\section{Finding Quiet Detail}\label{sec:detail}

Separate an acquisition change from a viewing change by recoloring the same
frozen history. Use a sound with both strong and weak partials or a long decay.

\paragraph{Starting point}
Load \textbf{DecibelInspection} and capture a few seconds of the sound.
Keep input gain at $0\,\mathrm{dB}$ and Slope at zero. Then press Run to freeze.

\begin{enumerate}
  \item Hover over a strong band and a quieter region. In Decibels mode,
  Raw reports the stored bin below the cursor frequency; Color reports the
  weighted, interpolated image sample. Neither readout is clamped to the
  palette's endpoints.
  \item Raise Floor from $-90$ to $-60\,\mathrm{dB}$ using its right-click
  numeric entry. Weak detail now shares the bottom color. This changes
  contrast; it does not remove noise or rewrite the stored spectra.
  \item Lower Floor again and lower Ceil toward the quiet detail. More of
  the palette is allocated to that range, while strong peaks can saturate at
  the top color. Raise Ceil again to distinguish the strong values.
  \item Try Cividis or Gray to compare color and lightness cues. Switch to
  Linear and back to Decibels. Each scale restores its own range; the stored
  history is the same throughout the experiment.
\end{enumerate}

\paragraph{A worked contrast example}
With Floor $-90$ and Ceil $0\,\mathrm{dB}$, a $-60\,\mathrm{dB}$ sample
sits one third of the way through the palette. With Floor $-70$ and Ceil
$-50\,\mathrm{dB}$, that same sample sits halfway through it. Values above
$-50\,\mathrm{dB}$ all reach the top color in the second view. Palette
position is numeric; equal steps need not look equally bright in every map.

\begin{center}
\renewcommand{\arraystretch}{1.2}
\begin{tabularx}{\textwidth}{@{}p{.35\textwidth}X@{}}
\toprule
\textbf{Change while frozen} & \textbf{Effect} \\
\midrule
Palette, range, Slope & Recolors retained history. \\
Freq Scale, LO/HI Freq & Remaps or crops the frequency view. \\
Window, Average, Smooth & Takes effect on new analysis after resuming. \\
Input gain, AC coupling & Affects samples captured after resuming. \\
\bottomrule
\end{tabularx}
\end{center}

To compare analysis settings, resume and allow fresh history to replace the
old columns. Save an external image to preserve a result: a saved Rack patch
retains settings and Run state, but not the captured history.
